\documentclass[aspectratio=169]{beamer}
\usepackage[utf8]{inputenc}
\usepackage[brazilian]{babel}
\usepackage{graphicx}
\usepackage{xcolor}

% UFG Colors (approximate standard blue/yellow)
\definecolor{ufgblue}{RGB}{0, 77, 133}
\definecolor{ufgyellow}{RGB}{255, 204, 0}

\mode<presentation> {
  \usetheme{Boadilla}
  \usecolortheme[named=ufgblue]{structure}
  \setbeamertemplate{navigation symbols}{}
  \setbeamertemplate{footline}[frame number]
  \setbeamercolor{title}{fg=ufgblue, bg=ufgyellow!20}
  \setbeamercolor{frametitle}{fg=white, bg=ufgblue}
  \setbeamercolor{itemize item}{fg=ufgyellow}
}

\title[Evolução e Refatoração Full-Stack do SIPROS]{Relato de Experiência e Pesquisa-Ação: Evolução e Refatoração Full-Stack do Sistema SIPROS}
\author[Phablo Tavares Paixão]{Phablo Tavares Paixão}
\institute[UFG]{
  Universidade Federal de Goiás - UFG\\
  Instituto de Informática\\
  Bacharelado em Engenharia de Software\\[1ex]
  \textit{Orientadora:} Profa. Dra. Sofia Larissa da Costa Paiva\\
  \textit{Coorientador:} Prof. Dr. Juliano Lopes de Oliveira
}
\date{2026}

\begin{document}

\begin{frame}
  \titlepage
\end{frame}

\begin{frame}{1. Introdução e Contexto}
  \begin{itemize}
    \setlength\itemsep{1em}
    \item \textbf{O Projeto:} Sistema SIPROS.
    \item \textbf{Onde:} Disciplina de Fábrica de Software (2026/1).
    \item \textbf{A Equipe:} "Equipe Dourada".
    \item \textbf{O Problema Original:} Base legada com dados \textit{mockados}, ausência de segurança real e falta de validação estruturada.
    \item \textbf{A Solução:} Refatoração \textit{full-stack} orientada por qualidade e segurança.
  \end{itemize}
\end{frame}

\begin{frame}{2. Objetivos e Metodologia}
  \begin{itemize}
    \setlength\itemsep{1em}
    \item \textbf{Objetivo Geral:} Evoluir o sistema de um protótipo \textit{mockado} para uma aplicação segura e validada.
    \item \textbf{Metodologia:} 
    \begin{itemize}
        \item Pesquisa-Ação: Intervenções contínuas na base de código seguidas de avaliação.
        \item Práticas Ágeis: Inspirado em Kanban e Scrum (flexibilizado para o contexto acadêmico).
    \end{itemize}
  \end{itemize}
\end{frame}

\begin{frame}{3. Foco em Segurança: OAuth 2.0 e PKCE}
  \begin{itemize}
    \setlength\itemsep{1em}
    \item \textbf{Desafio:} SPAs (Single Page Applications) não podem armazenar segredos de cliente de forma segura.
    \item \textbf{Implementação:}
    \begin{itemize}
        \item Integração com \textbf{Keycloak} para gestão de identidades.
        \item Adoção do fluxo \textbf{OAuth 2.0 com PKCE} (Proof Key for Code Exchange).
        \item Login social via Google habilitado de forma segura no \textit{front-end} Angular.
    \end{itemize}
  \end{itemize}
\end{frame}

\begin{frame}{4. Qualidade e Verificação \& Validação (V\&V)}
  \begin{itemize}
    \setlength\itemsep{1em}
    \item \textbf{Alinhamento de Requisitos:} Mapeamento rigoroso de Casos de Uso versus interfaces no Figma.
    \item \textbf{Garantia de Qualidade (QA):}
    \begin{itemize}
        \item Implementação de políticas rigorosas de \textbf{Code Review} via GitHub (Pull Requests).
        \item Inserção de \textbf{Testes de Unidade} automatizados.
        \item Foco nas métricas de qualidade de software (ISO 25010).
    \end{itemize}
  \end{itemize}
\end{frame}

\begin{frame}{5. Resultados Alcançados}
  \begin{itemize}
    \setlength\itemsep{1em}
    \item Autenticação e autorização funcionais no ambiente de Fábrica.
    \item Fluxo de desenvolvimento estabelecido (rastreabilidade desde o requisito até o PR aprovado).
    \item Código mais maduro, testável e manutenível em comparação à base legada inicial.
  \end{itemize}
\end{frame}

\begin{frame}{6. Lições Aprendidas}
  \begin{itemize}
    \setlength\itemsep{1em}
    \item \textbf{Hard Skills:} Arquitetura de segurança, interceptadores HTTP, testes no ecossistema Angular.
    \item \textbf{Soft Skills:} Trabalho colaborativo, resolução de conflitos, resiliência perante o código legado.
    \item \textbf{Conselhos aos futuros alunos:}
    \begin{itemize}
        \item Abrace o código legado como oportunidade de aprendizado.
        \item Construa a cultura de testes desde o \textit{Sprint} 1.
    \end{itemize}
  \end{itemize}
\end{frame}

\begin{frame}{7. Trabalhos Futuros e Conclusão}
  \begin{itemize}
    \setlength\itemsep{1em}
    \item \textbf{Trabalhos Futuros:}
    \begin{itemize}
        \item Alcançar 100\% de cobertura de testes no \textit{back-end}.
        \item Implementação de testes \textit{End-to-End} (E2E).
    \end{itemize}
    \vspace{0.5em}
    \item \textbf{Conclusão:} A Fábrica de Software não apenas entregou valor técnico ao SIPROS, mas funcionou como um simulador fidedigno dos desafios da indústria real de engenharia de software.
  \end{itemize}
\end{frame}

\begin{frame}[plain, c]
  \begin{center}
    \Huge \textbf{Obrigado(a)!}
    
    \vspace{1em}
    \Large Phablo Tavares Paixão
  \end{center}
\end{frame}

\end{document}
